\documentclass[11pt,a4paper]{article}
\usepackage[margin=24mm]{geometry}
\usepackage[T1]{fontenc}
\usepackage{lmodern,microtype,amsmath,amssymb,amsthm,mathtools}
\usepackage{booktabs,array,enumitem}
\usepackage[hidelinks]{hyperref}
\setlength{\parindent}{0pt}
\setlength{\parskip}{5pt}
\newtheorem{theorem}{Theorem}[section]
\newtheorem{proposition}[theorem]{Proposition}
\newtheorem{corollary}[theorem]{Corollary}
\newcommand{\Q}{\mathbb Q}
\newcommand{\E}{\mathbb E}
\newcommand{\F}{\mathcal F}
\newcommand{\dd}{\,\mathrm d}
\DeclareMathOperator{\Var}{Var}

\title{The exact DLV--Black product lift\\
from the Wiese--Wissel drift condition to a smooth SANOS surface}
\author{A self-contained construction and proof}
\date{24 September 2026}

\begin{document}
\maketitle

\begin{abstract}
This note proves the particular construction of interest.  The factor $X$ is
the discrete martingale encoded by a dynamically arbitrage-free DLV density
tower.  Its DLV coordinates obey the Wiese--Wissel drift condition: Wiese's
density-coordinate form and Wissel's price-coordinate form are algebraically
equivalent.
An independent positive Black--Scholes martingale $Y$ supplies the smoothing.
For $S=XY$, the conditional call surface is a finite SANOS mixture.  An exact
It\^o calculation proves zero drift not only at the inherited DLV nodes but at
every strike between them.  The proof also identifies the necessary
orthogonality, maturity-clock and grid-recalibration qualifications.
\end{abstract}

\section{What is being identified}

The roles of the two factors are different:
\begin{equation}
 \boxed{
 \begin{array}{ll}
 X:&\text{the discrete martingale factor generated by the DLV density tower},\\[1mm]
 Y:&\text{an independent positive Black martingale generating the smooth kernel},\\[1mm]
 S=XY:&\text{the observable smoothed underlying.}
 \end{array}}
 \label{eq:roles}
\end{equation}
Strictly speaking, $X$ is not ``the DLV drift''.  The DLV model determines the
conditional laws of $X$ and the dynamics of the DLV coordinates.  The
Wiese--Wissel condition determines the risk-neutral drift of the local
variance coordinate $\ell$.  The stock factor $X$ itself is a martingale.

We work in discounted units.  Let
\[
 0<x_0<\cdots<x_N
\]
be the latent DLV support and let $T_1<\cdots<T_M$ be the model expiries.  A
coherent DLV tower has conditional masses
\begin{equation}
 q_i^m(t)=\Q(X_{T_m}=x_i\mid\F_t^X),
 \qquad t\leq T_m,
 \label{eq:qconditional}
\end{equation}
so that
\begin{equation}
 q_i^m\geq0,\qquad \sum_iq_i^m=1,\qquad
 \sum_i x_iq_i^m=X_t .
 \label{eq:massmean}
\end{equation}
For fixed $(m,i)$, $q_i^m(t)$ is a bounded martingale.  In the continuous case
write
\begin{equation}
 \dd q_i^m(t)=\psi_i^m(t)\cdot\dd W_t^X,
 \qquad
 \dd X_t=\sum_i x_i\psi_i^m(t)\cdot\dd W_t^X.
 \label{eq:qdynamics}
\end{equation}
The second formula is independent of $m$ by the common-forward condition.

The hypothesis that the DLV tower is coherent is material.  It is supplied,
for example, by a DLV/Wissel model satisfying its existence assumptions, or by
finite martingale transition kernels between the expiry marginals.  A list of
unrelated static density fits is not yet such a process.

\section{The \texorpdfstring{$X$}{X}-only DLV drift condition}

Define the $X$-only calls on a fixed cash strike grid $K_n$ by
\begin{equation}
 c_n^m(t)=\E[(X_{T_m}-K_n)^+\mid\F_t^X]
          =\sum_iq_i^m(t)(x_i-K_n)^+ .
 \label{eq:xcall}
\end{equation}
Every $c_n^m$ is a martingale.  Its discrete butterfly is
\begin{equation}
 \widehat q_n^m=
 \frac{c_{n+1}^m-c_n^m}{K_{n+1}-K_n}
 -\frac{c_n^m-c_{n-1}^m}{K_n-K_{n-1}} .
 \label{eq:butterfly}
\end{equation}
When the strike and support grids coincide, with the usual boundary masses,
$\widehat q^m$ is precisely the conditional mass vector in
\eqref{eq:qconditional}.  We retain a hat in this section to show which fact is
used: $\widehat q_n^m$ is a deterministic linear combination of traded calls
and therefore has zero drift.

For a back-month slab define $\ell_n^m$ by the exact discrete Dupire identity
\begin{equation}
 c_n^m-c_n^{m-1}
 =\omega_n^m\ell_n^m\widehat q_n^m,
 \qquad
 \omega_n^m=
 \frac{K_n^2(T_m-T_{m-1})}{K_{n+1}-K_{n-1}}.
 \label{eq:dlvidentity}
\end{equation}
Assume $\widehat q_n^m>0$ and write
\begin{equation}
 \dd\widehat q_n^m=\psi_n^m\cdot\dd W_t^X,
 \qquad
 \dd\ell_n^m=\alpha_n^m\dd t+\beta_n^m\cdot\dd W_t^X.
 \label{eq:ellq}
\end{equation}

\begin{proposition}[Wiese--Wissel drift condition]
For every back-month node,
\begin{equation}
 \boxed{
 \alpha_n^m=-\frac{\beta_n^m\cdot\psi_n^m}{\widehat q_n^m}
            =-\eta_n^m\cdot\beta_n^m,
 \qquad \eta_n^m:=\frac{\psi_n^m}{\widehat q_n^m}.}
 \label{eq:magnus}
\end{equation}
Under the exact identification $\ell_n^m=(x_n^m)^2$, this is the
price-coordinate form of the Wiese--Wissel risk-neutral drift condition.
\end{proposition}

\begin{proof}
The left-hand side of \eqref{eq:dlvidentity} is the difference of two
martingales.  The drift of its right-hand side is
\[
 \omega_n^m\bigl(widehat q_n^m\alpha_n^m
                 +\beta_n^m\cdot\psi_n^m\bigr).
\]
It must vanish.  Division by $\omega_n^m\widehat q_n^m$ proves
\eqref{eq:magnus}.

For the price-space dictionary, suppose Wissel writes
\[
 \frac{\dd x_n^m}{x_n^m}
 =\bar u_n^m\dd t+v_n^m\cdot\dd W_t^X .
\]
Since $\ell=(x)^2$, It\^o's formula gives
\[
 \beta_n^m=2\ell_n^mv_n^m,
 \qquad
 \frac{\alpha_n^m}{\ell_n^m}
 =2\bar u_n^m+\lVert v_n^m\rVert^2.
\]
The relative calendar-spread volatility is
\[
 R_n^m=\frac{\beta_n^m}{\ell_n^m}
       +\frac{\psi_n^m}{\widehat q_n^m}
       =2v_n^m+\eta_n^m.
\]
Substitution in Wissel's risk-neutral condition
$\bar u_n^m=(-R_n^m+\tfrac32v_n^m)\cdot v_n^m$ gives
\[
 \frac{\alpha_n^m}{\ell_n^m}
 =-2\eta_n^m\cdot v_n^m
 =-\eta_n^m\cdot\frac{\beta_n^m}{\ell_n^m},
\]
which is exactly \eqref{eq:magnus}.
\end{proof}

Thus Wiese's density-coordinate derivation and Wissel's price-coordinate
derivation impose the same grid drift.  We call their common restriction the
\emph{Wiese--Wissel drift condition}.  This is the drift restriction inherited
by the $X$ factor in the product construction.

With jumps, the exact predictable-drift condition is
\begin{equation}
 \alpha_n^m=-\frac{\beta_n^m\cdot\psi_n^m
 +\displaystyle\int\Delta\ell_n^m(z)\Delta\widehat q_n^m(z)\,
                    \nu_t(\dd z)}{\widehat q_n^m}.
 \label{eq:magnusjump}
\end{equation}
The continuous formula is sufficient for the smooth lift proved next.

\section{The independent Black smoothing factor}

Let $W^Y$ be Brownian motion independent of $\F^X_\infty$, and set
\begin{equation}
 \frac{\dd Y_t}{Y_t}=\sigma_Y(t)\cdot\dd W_t^Y,
 \qquad
 Y_0=1,
 \label{eq:Y}
\end{equation}
where $\sigma_Y$ is deterministic and square-integrable.  Put
\begin{equation}
 V(t)=\int_0^t\lVert\sigma_Y(u)\rVert^2\dd u,
 \qquad
 w(t,T)=V(T)-V(t).
 \label{eq:clock}
\end{equation}
Then $Y$ is a positive true martingale and, conditionally on $\F_t^Y$,
\[
 \frac{Y_T}{Y_t}
 =\exp\!\left(-\frac12w(t,T)+\sqrt{w(t,T)}Z\right),
 \qquad Z\sim N(0,1).
\]

Let
\begin{equation}
 B(a,K,w)=\E[(ae^{-w/2+\sqrt w Z}-K)^+]
 \label{eq:black}
\end{equation}
denote the Black call kernel.  It satisfies the heat identity
\begin{equation}
 B_w(a,K,w)=\frac12a^2B_{aa}(a,K,w).
 \label{eq:heat}
\end{equation}
The function $B$ is the smoothing operator; $Y$ is the martingale whose
conditional transition law produces that operator.

\section{Exact product-lift identity}

\begin{proposition}[DLV--Black product lift]
\label{thm:product}
Let $X$ be a positive, bounded DLV martingale with coherent conditional
density tower \eqref{eq:qconditional}, whose back-month coordinates satisfy
\eqref{eq:magnus}.  Let $Y$ be the independent Black martingale
\eqref{eq:Y}, and define
\begin{equation}
 S_t=X_tY_t.
 \label{eq:S}
\end{equation}
Then $S$ is a positive true martingale.  For every model expiry $T_m$, every
trading time $t\leq T_m$, and every cash strike $K>0$, its call price is
\begin{equation}
 \boxed{
 C_t(T_m,K)=
 \sum_{i=0}^{N}q_i^m(t)
 B\!\left(x_iY_t,K,w(t,T_m)\right).}
 \label{eq:sanos}
\end{equation}
Consequently:
\begin{enumerate}[label=(\roman*)]
\item \eqref{eq:sanos} is a finite SANOS representation with weights $q_i^m$,
anchors $A_i(t)=x_iY_t$, and common smoothing variance $w(t,T_m)$;
\item it is statically arbitrage-free in strike and maturity;
\item for every fixed $(T_m,K)$, including every $K$ strictly between DLV
nodes, $C_t(T_m,K)$ is a true $\Q$-martingale;
\item at expiry it settles to the ordinary payoff $(S_{T_m}-K)^+$.
\end{enumerate}
\end{proposition}

\begin{proof}
Independence and positivity give
\[
 \E[S_T\mid\F_t^X\vee\F_t^Y]
 =\E[X_T\mid\F_t^X]\E[Y_T\mid\F_t^Y]=X_tY_t=S_t,
\]
so $S$ is a true martingale.  Condition first on $X_{T_m}=x_i$ and then on
$\F_t^X\vee\F_t^Y$.  The conditional law of $Y_{T_m}/Y_t$ is lognormal with
total variance $w(t,T_m)$, hence
\[
 \E[(S_{T_m}-K)^+\mid\F_t]
 =\sum_iq_i^m(t)B(x_iY_t,K,w(t,T_m)),
\]
which proves \eqref{eq:sanos}.  It is a conditional expectation of an
ordinary call payoff on a single martingale.  Dynamic no-arbitrage and
calendar monotonicity follow.  Positivity of the weights and the Black kernels
give monotonicity and convexity in strike.

At $t=T_m$, $w(T_m,T_m)=0$, and the conditional mass collapses onto the
realised state $X_{T_m}=x_I$.  Therefore
\[
 C_{T_m}(T_m,K)=B(x_IY_{T_m},K,0)
 =(S_{T_m}-K)^+.
\]
\end{proof}

The proposition gives the exact construction needed: the DLV factor supplies
the state-dependent mixture weights and their constrained dynamics; the Black
factor supplies the common smoothing; their product supplies one underlying
whose entire smooth call surface is dynamically arbitrage-free.

\section{The between-grid drift condition, calculated explicitly}

The conditional-expectation proof is decisive, but the generator calculation
shows exactly how the grid drift is extended.  Fix $(T_m,K)$ and abbreviate
\[
 a_i(t)=x_iY_t,\qquad w_m(t)=w(t,T_m),\qquad
 G_i(t)=B(a_i(t),K,w_m(t)).
\]
Since
\[
 \dd a_i=a_i\sigma_Y\cdot\dd W^Y,\qquad
 \dd w_m=-\lVert\sigma_Y\rVert^2\dd t,
\]
It\^o's formula and \eqref{eq:heat} give
\begin{align}
 \dd G_i={}&B_a(a_i,K,w_m)a_i\sigma_Y\cdot\dd W^Y\notag\\
 &+\left[-\lVert\sigma_Y\rVert^2B_w
 +\frac12\lVert\sigma_Y\rVert^2a_i^2B_{aa}\right]\dd t\notag\\
 ={}&B_a(a_i,K,w_m)a_i\sigma_Y\cdot\dd W^Y.
 \label{eq:kernelmart}
\end{align}
Thus every moving Black kernel is itself a martingale.  Now
\begin{equation}
 C_t(T_m,K)=\sum_iq_i^m(t)G_i(t).
\end{equation}
Integration by parts gives
\begin{align}
 \dd C={}&sum_iG_i\,\dd q_i^m
          +\sum_iq_i^m\,\dd G_i
          +\sum_i\dd\langle q_i^m,G_i\rangle_t.
 \label{eq:productito}
\end{align}
The first two sums have zero drift.  Independence gives
$\dd\langle W^X,W^Y\rangle_t=0$, so the last sum vanishes.  Hence
\begin{equation}
 \boxed{
 \operatorname{drift}\,C_t(T_m,K)=0
 \quad\text{for every }K>0.}
 \label{eq:offgridzero}
\end{equation}

Equivalently, the full off-grid It\^o restriction is
\begin{equation}
 \boxed{
 \sum_i B_i\,\mu_{q_i}
 +\sum_iq_i\left(
 -\lVert\sigma_Y\rVert^2B_{w,i}
 +\tfrac12\lVert\sigma_Y\rVert^2a_i^2B_{aa,i}
 \right)
 +\sum_i a_iB_{a,i},
       \frac{\dd\langle q_i,\log Y\rangle_t}{\dd t}=0.}
 \label{eq:extendeddrift}
\end{equation}
In the exact construction, its three lines vanish respectively because
\begin{equation}
 \mu_{q_i}=0,\qquad
 B_w=\tfrac12a^2B_{aa},\qquad
 \dd\langle q_i,\log Y\rangle=0.
 \label{eq:threecancel}
\end{equation}
Equation \eqref{eq:extendeddrift} is the desired between-node extension of the
finite DLV drift condition.  It is valid for every strike, not only for a
dense numerical audit mesh.

If $W^X$ and $W^Y$ are correlated, the last term in
\eqref{eq:extendeddrift} is generally nonzero.  The simple product lift then
requires an additional compensating drift.  Independence, or at least
orthogonality of the relevant martingale parts, is therefore a structural
condition rather than a cosmetic assumption.

\section{How the node condition and the smooth condition fit together}

There are two exact but distinct statements.

First, on the latent $X$ grid, the DLV identity
\[
 c_n^m-c_n^{m-1}
 =\omega_n^m\ell_n^m\widehat q_n^m
\]
and its product rule give the Wiese--Wissel drift condition
\[
 \alpha_n^m=-\eta_n^m\cdot\beta_n^m.
\]

Second, after multiplication by $Y$, the actual smoothed calls satisfy
\[
 C_t(T_m,K)=\sum_iq_i^mB(x_iY_t,K,w_m(t))
\]
and the extended drift identity \eqref{eq:extendeddrift} vanishes for every
strike.  If one samples these actual smooth calls on any strike grid and
recomputes their butterflies and DLV coordinates, those newly computed
coordinates also satisfy the Wiese--Wissel drift condition, because they are
deterministic functions of martingale calls.  They need not be numerically
equal to the latent $X$-only DLV coordinates.

This last point is essential.  Multiplying an already calibrated $X$ by a
nontrivial $Y$ generally changes the prices at the original cash strikes:
\begin{equation}
 \sum_iq_i^m B(x_iY_t,K_n,w_m(t))
 \neq
 \sum_iq_i^m(x_i-K_n)^+ .
 \label{eq:notpreserve}
\end{equation}
Therefore the product theorem is an exact dynamic construction, but it is not
an automatic price-preserving post-processing of an arbitrary pre-existing
DLV grid.

To match a prescribed finite call grid $C_n^{m,\mathrm{target}}$ exactly, the
initial product model must be calibrated through the linear constraints
\begin{equation}
 \sum_iq_i^m B(x_iY_0,K_n,w(0,T_m))
 =C_n^{m,\mathrm{target}},\qquad
 \sum_iq_i^m=1,\qquad
 \sum_iq_i^mx_i=X_0,\qquad q_i^m\geq0,
 \label{eq:inverse}
\end{equation}
together with martingale-transport consistency across $m$.  If these
constraints are feasible, the theorem supplies an exact dynamically
arbitrage-free smooth extension of that grid.  If they are infeasible for the
chosen anchors and smoothing clock, one must move or enlarge the anchors,
reduce the smoothing, or fit within bid--ask.  The drift proof cannot repair a
failed static inverse problem.

\section{Front expiry and the removal of observed-spot pinning}

The construction has a particularly clean front-expiry interpretation.  The
latent factor $X_{T_1}$ lands on the finite DLV support, but the observed spot
is
\[
 S_{T_1}=X_{T_1}Y_{T_1}.
\]
For a nondegenerate Black factor, $S_{T_1}$ has continuous support and is not
pinned to the fixed cash strike grid.  Nevertheless the call settles exactly,
because $w(t,T_1)\downarrow0$ and the surviving moving anchor is
$X_{T_1}Y_{T_1}=S_{T_1}$.

Thus the product construction separates two questions:
\begin{enumerate}
\item the DLV factor $X$ carries the finite martingale transport and obeys the
Wiese--Wissel drift condition;
\item the Black factor $Y$ turns its atoms into a smooth observable call
distribution while preserving one common underlying and ordinary settlement.
\end{enumerate}
It is an exact front-expiry construction, not merely a terminal interpolation.
Its empirical adequacy and the feasibility of the inverse calibration
\eqref{eq:inverse} are separate numerical questions.

\section{Conclusion}

The exact theorem can be summarised as
\begin{equation}
 \boxed{
 \begin{gathered}
 \text{coherent DLV martingale }X
 \quad+\quad
 \text{independent Black martingale }Y\\
 \Longrightarrow\quad
 S=XY,\qquad
 C_t(T,K)=\sum_iq_i(t,T)B(x_iY_t,K,V(T)-V(t)),\\
 \operatorname{drift}C_t(T,K)=0
 \quad\text{for every fixed }(T,K).
 \end{gathered}}
\end{equation}
On the $X$ grid, dynamic no-arbitrage is encoded by
$\alpha=-\eta\cdot\beta$.  Between grid points, it is encoded by the Black heat
equation together with orthogonality of the $X$ and $Y$ martingale parts.  The
two cancellations are not competing formulas: they are the grid and smooth
parts of the same product construction.

\begin{thebibliography}{9}
\bibitem{Wiese}
M. Wiese,
\emph{A drift restriction for discrete local volatility: the argument in five
steps}, discussion note, 2026.

\bibitem{SchweizerWissel}
M. Schweizer and J. Wissel,
\emph{Arbitrage-free market models for option prices: the multi-strike case},
Finance and Stochastics 12 (2008), 469--505.

\bibitem{BuehlerRyskin}
H. Buehler and E. Ryskin,
\emph{Discrete local volatility for large time steps}, SSRN 2642630.

\bibitem{SANOS}
H. Buehler, B. Horvath, A. Kratsios, Y. Limmer and R. Saqur,
\emph{SANOS: Smooth strictly Arbitrage-free Non-parametric Option Surfaces},
arXiv:2601.11209.
\end{thebibliography}

\end{document}
